\documentclass[11pt]{article}
\usepackage[margin=1in]{geometry}
\usepackage{amsmath,amssymb,amsthm}
\usepackage[colorlinks=true,linkcolor=blue,citecolor=blue,urlcolor=blue]{hyperref}

\newtheorem{theorem}{Theorem}
\newtheorem{lemma}{Lemma}
\newtheorem{corollary}{Corollary}
\newtheorem{proposition}{Proposition}
\theoremstyle{definition}
\newtheorem{definition}{Definition}
\newtheorem{remark}{Remark}
\newtheorem{question}{Question}

\newcommand{\Z}{\mathbb{Z}}
\newcommand{\abs}[1]{\left\lvert #1 \right\rvert}

\title{A Balanced Covering Pair of Size $\sqrt{n}$ and Polynomial Magnitude:\\
Resolving a Difference-Set Question for the Smoothness Covering Decomposition}
\author{Eric Schultz \qquad Claude (Anthropic)}
\date{July 2026}

\begin{document}
\maketitle

\begin{abstract}
A companion note to the Smoothness Covering Decomposition showed how to cover $[1,n]$ by
divisibility with a set $D$ of size $\pi(n)(1+o(1))$ and elements of magnitude
$\exp(O(\sqrt n))$, and asked (Remark~3 there) for the least $K$ such that a
\emph{balanced} pair $(S,T)$ with $\abs{S}=\abs{T}=K$ covers $[1,n]$ by nonzero
differences --- every $i\le n$ dividing some $s-t$ with $s\in S$, $t\in T$, $s\ne t$. That
remark cautioned the balanced realization looked infeasible for $K\ll\abs{D}$. We show the
opposite: an elementary two-scale grid gives a balanced pair with
$K=\lceil\sqrt{n+1}\,\rceil=(1+o(1))\sqrt n$ and all elements of magnitude at most
$n+O(\sqrt n)$. This improves the trivial realization simultaneously on \emph{both} axes
--- size $\sqrt n$ versus $\pi(n)$, and magnitude $O(n)$ versus $\exp(O(\sqrt n))$. We also
observe that minimizing $K$ alone is degenerate ($K=1$ via a single least-common-multiple
element, at magnitude $\exp(O(n))$), so the meaningful object is the size--magnitude
trade-off; the grid occupies one clean corner of it, and the shape of the full curve
(in particular, any nontrivial lower bound on $K$ at polynomial magnitude) is left open.
The finite claims were verified by direct enumeration. This note has not been externally
refereed.
\end{abstract}

\section{Setup}

We work with the divisibility-covering notion of the companion decomposition.

\begin{definition}
A pair of finite sets $S,T\subset\Z$ is a \emph{covering pair} for $[1,n]$ if for every
integer $i$ with $1\le i\le n$ there exist $s\in S$, $t\in T$ with $s\ne t$ and
$i\mid(s-t)$. It is \emph{balanced} if $\abs S=\abs T$; we write $K=\abs S=\abs T$. Its
\emph{magnitude} is $\max\{\abs x : x\in S\cup T\}$.
\end{definition}

The nonzero requirement $s\ne t$ is essential: without it the difference $0$ is divisible
by every $i$, and $S=T=\{0\}$ would ``cover'' $[1,n]$ vacuously. With it, the problem is
the genuine one raised in Remark~3 of the companion note.

The trivial realization $S=D$, $T=\{0\}$ from a covering set $D$ gives $K=\abs D$; for the
decomposition's $D$ that is $\pi(n)(1+o(1))$ at magnitude $\exp(O(\sqrt n))$. Remark~3
there worried that a balanced realization with $K\ll\abs D$ is obstructed, because
prescribing the $\abs D$ specific values of $D$ as differences imposes $\abs D$ constraints
on the $2K$ values of $S\cup T$. The resolution below is to abandon reproducing $D$'s
specific values: a balanced pair need not realize $D$ at all, only cover $[1,n]$ afresh.

\section{The construction}

\begin{theorem}[Balanced covering pair]\label{thm:main}
For every $n\ge1$ let $m=\lceil\sqrt{n+1}\,\rceil$ and set
\[
  S=\{0,1,2,\dots,m-1\},\qquad T=\{0,-m,-2m,\dots,-(m-1)m\}.
\]
Then $(S,T)$ is a balanced covering pair for $[1,n]$ with $K=m=(1+o(1))\sqrt n$ and
magnitude $(m-1)m< n+O(\sqrt n)$.
\end{theorem}

\begin{proof}
For $s=a\in S$ and $t=-bm\in T$ ($0\le a,b\le m-1$) the difference is
$s-t=a+bm$. As $a$ ranges over $\{0,\dots,m-1\}$ and $b$ over $\{0,\dots,m-1\}$, the value
$a+bm$ takes \emph{every} integer in $\{0,1,\dots,m^2-1\}$ exactly once (base-$m$
representation). Since $m=\lceil\sqrt{n+1}\,\rceil$ gives $m^2\ge n+1$, the difference set
$\{s-t\}$ contains $\{1,2,\dots,n\}$. Hence for each $i\le n$ the value $i$ itself is a
nonzero difference, and trivially $i\mid i$; the pair covers $[1,n]$. Both sets have $m$
elements, so $K=m=(1+o(1))\sqrt n$. The largest magnitude is $\abs{-(m-1)m}=(m-1)m$, and
with $m<\sqrt{n+1}+1$ this is below $n+O(\sqrt n)$.
\end{proof}

\begin{remark}[Two-axis improvement]
The pair of Theorem~\ref{thm:main} improves the trivial realization on both axes at once:
size $\sqrt n$ against $\pi(n)\sim n/\ln n$, and magnitude $O(n)$ against the
decomposition's $\exp(O(\sqrt n))$. It is therefore not merely a balanced substitute for
$D$ but a strictly cheaper covering object in both parameters. The obstruction of Remark~3
in the companion note does not bind because the construction never reproduces the specific
elements of $D$; it covers $[1,n]$ directly through the $m^2$ differences of an
$m\times m$ grid.
\end{remark}

\section{Minimizing $K$ alone is degenerate; the trade-off is the real object}

\begin{proposition}\label{prop:lcm}
For every $n$ the pair $S=\{0\}$, $T=\{\operatorname{lcm}(1,\dots,n)\}$ is a balanced
covering pair for $[1,n]$ with $K=1$. Its magnitude is
$\operatorname{lcm}(1,\dots,n)=\exp((1+o(1))n)$.
\end{proposition}

\begin{proof}
The single nonzero difference is $-\operatorname{lcm}(1,\dots,n)$, which every $i\le n$
divides. The magnitude estimate is the prime number theorem for the second Chebyshev
function, $\log\operatorname{lcm}(1,\dots,n)=\psi(n)\sim n$.
\end{proof}

Thus the bare question ``least $K$'' has answer $1$ and is not the interesting one: the
cost has simply migrated into magnitude, from $\exp(O(\sqrt n))$ up to $\exp(O(n))$. The
meaningful object is the \emph{size--magnitude trade-off}: the least $K$ achievable subject
to a magnitude bound $M$. Three points on this curve are now known:
\[
  \begin{array}{lll}
    \text{(i)}   & K=1,          & M=\exp(O(n))       \quad\text{(Proposition~\ref{prop:lcm})};\\
    \text{(ii)}  & K=\pi(n),     & M=\exp(O(\sqrt n)) \quad\text{(decomposition's }D\text{)};\\
    \text{(iii)} & K=\sqrt n,    & M=O(n)             \quad\text{(Theorem~\ref{thm:main})}.
  \end{array}
\]
Point (iii) dominates (ii) in both coordinates, so (ii) is not on the lower boundary of the
trade-off. The genuinely open direction is the boundary near (iii):

\begin{question}
At polynomial magnitude $M=n^{O(1)}$, is $K=\Theta(\sqrt n)$ optimal, or can a balanced
covering pair have $K=o(\sqrt n)$? Equivalently, is there a lower bound $K=\Omega(\sqrt n)$
for covering pairs of polynomially bounded magnitude?
\end{question}

A lower bound here cannot come from a crude counting of differences alone: covering needs
only a \emph{multiple} of each $i\le n$ among the $\le K^2$ differences, not $i$ itself, and
the relaxation to multiples is what makes small-$n$ optima dip below $\lceil\sqrt n\rceil$
(see Section~\ref{sec:verify}). Whether that relaxation buys an asymptotic improvement at
polynomial magnitude, or is washed out, is what the question asks.

\section{Verification}\label{sec:verify}

The construction of Theorem~\ref{thm:main} was checked by direct enumeration: for
$n\in\{16,100,1000,10^4\}$ the pair $(S,T)$ leaves no $i\le n$ uncovered, and the measured
$K/\sqrt n$ is $1.25,1.10,1.01,1.01$ respectively, converging to $1$. The magnitudes
$(m-1)m$ are $20,110,992,10100$ against $n=16,100,1000,10^4$, confirming the $O(n)$ bound.
Proposition~\ref{prop:lcm} was checked for $n\le500$: the single element
$\operatorname{lcm}(1,\dots,n)$ is divisible by every $i\le n$.

Exhaustive minimum-$K$ search (nonzero differences, values bounded by $2n$) shows the
multiple-relaxation genuinely dips below $\lceil\sqrt n\rceil$ at small $n$: the least $K$
is $2$ up to $n=10$ and $3$ up to $n=20$, whereas $\lceil\sqrt n\rceil$ reaches $3$ already
at $n=5$ and $4$ at $n=10$. For instance $S=\{0,2\},T=\{16,20\}$ covers $[1,10]$ with $K=2$
(differences $\{-14,-16,-18,-20\}$, whose divisors include every $i\le10$). These are the
data motivating the open question; they do not settle the asymptotic rate, since the search
is exponential and was carried out only for $n\le24$.

\section*{Acknowledgement of scope}
The upper bound of Theorem~\ref{thm:main} is complete and elementary. The lower-bound side
of the trade-off --- and hence the exact optimal $K$ at polynomial magnitude --- is not
resolved here.

\begin{thebibliography}{9}
\bibitem{smooth} E. Schultz, \emph{A Smoothness Covering Decomposition: Covering $[1,n]$
by Divisibility with One Dedicated Element per Large Prime Power}, 2026 (companion note;
Remark~3 poses the balanced-pair question resolved here).
\end{thebibliography}

\end{document}
